% Native reconstruction of Springer SGA4 tome I backmatter.
% Source SHA-256 f4f025b21f5fe3c2ef6a94ef525ae510e58b26ecade9712821f8e091df20057b.
% Physical pages 534--539 (zero-based533--538), printed pages520--525.
\chapter*{术语索引}
\addcontentsline{toc}{chapter}{术语索引}
\begingroup
\small
\setlength{\LTleft}{0pt}
\setlength{\LTright}{0pt}
\begin{longtable}{@{}p{.68\linewidth}p{.27\linewidth}@{}}
% SOURCE physical page534, printed520, visually checked against page-534.png.
可及 (函子) & I 9.2\\
可及 (对象) & I 9.3\\
子拓扑斯的闭包 & IV 9.4.8\\
Artin 的 (集合) & I 0, I 11.6\\
Artin 的 (对象) & I 8.12.6\\[.7em]
双覆盖 (态射、族) & II 5.2\\
笛卡儿、余笛卡儿 & I 10.1\\
$\underline U$-范畴 & I 1.1\\
滤过范畴 & I 2.7\\
(拓扑斯的) 点范畴 & IV 4.6.3\\
伪滤过范畴 & I 2.7\\
共尾 (函子、子范畴) & I 8.1.1\\
余生成 (子范畴) & I 7.1\\
比较 (引理) & III 4.1\\
常值、本质常值 (ind-对象) & I 8.4\\
覆盖 (筛) & II 1.1.1\\
覆盖 (态射、族) & II 5.2\\
覆盖 (族) & II 1.2\\
筛 & I 4.1\\[.7em]
支配 (拓扑斯态射) & IV 8.8\\
普遍有效满态射 (族) & II 2.5\\
满态射 (族) & I 10.3\\
普遍严格满态射 (筛、族) & II 2.5\\
普遍满态射 (族) & I 10.3\\
满态射、严格满态射 & I 10.3\\
有效满态射 & I 10.3\\
本质 (态射、点) & IV 7.6\\
展拓扑斯 (étendue) & IV 9.8.2\\
子拓扑斯的外部 & IV 9.4.8\\[.7em]
层 (取值于范畴 $D$) & II 6.1\\
层 (集合) & II 2.1\\
层 (态射) & IV 10.2\\
层 (粘合) & IV 9.6\\
基数滤过 & I 9.12\\
余连续函子 & III 2.1\\
% SOURCE physical page535, printed521, visually checked against page-535.png.
保守函子 & I 6.1\\
连续函子 & III 1.1\\
粘合函子 & IV 9.3.5\\
左正合 (右正合、正合) 函子 & I 2.4\\
纤维函子 & IV 4.6.3\\
忠实函子 & I 6.1\\
延零函子 & IV 11.3.1\\
标量限制函子 & IV 11.1.4\\
截面函子 & IV 4.3.6\\
纤维函子 & IV 6.0\\
子拓扑斯的边界 & IV 9.4.8\\[.7em]
生成、余生成 (族) & I 7.1\\
生成、余生成 (子范畴) & I 7.1, I 7.9\\
生成 (景的生成族) & II 3.0.1\\
生成 (景的拓扑生成族) & II 3.0.1\\
泛化 (点的) & IV 4.2.2\\
大拓扑斯 (拓扑空间的) & IV 4.10, IV 2.5\\[.7em]
模的正像 & IV 13.2.1\\
拓扑斯态射的像 & IV 9.1.7\\
逆像 (子拓扑斯的) & IV 9.1.6\\
逆像 (诱导拓扑斯的) & IV 5.10\\
模的逆像 & IV 13.2.2\\
不可达 (基数) & I 0\\
不可达 (基数) & I 11.5\\
包含 (态射) & IV 5.2\\
ind-伴随 & I 8.11.1\\
ind-对象 & I 8\\
严格 ind-对象 & I 8.12.1\\
ind-可表 (函子) & I 8.2\\
诱导 (拓扑斯) & IV 5.2\\
子拓扑斯的内部 & IV 9.4.8\\[.7em]
归纳极限 & I 2.3.1\\
普遍归纳极限 & I 2.5\\
投射极限 & I 2.1\\
$\underline U$-投射极限 & I 2.2.1\\
有限投射极限与归纳极限 & I 2.3.1\\
局部 (性质关系) & IV 8.2\\
局部化 (函子) & IV 5.4\\
局部化 (态射) & IV 5.2\\
% SOURCE physical page536, printed522, visually checked against page-536.png.
单态射性、单态射 & I 10.4\\
拓扑斯态射 & IV 3.1\\
态射 (景的) & IV 4.9\\
赋环拓扑斯态射 & IV 13.1\\
局部赋环拓扑斯态射 & IV 13.9\\[.7em]
始对象 & II 4.5\\
开集 (拓扑斯的) & IV 8.2\\[.7em]
小 (集合、群、环、范畴\ldots) & I 1.0\\
$\underline U$-小 (集合) & I 1.0\\
嵌入 & IV 9.1.1\\
闭嵌入 & IV 9.3.5\\
开嵌入 & IV 9.2.1\\
几何点 & IV 13.9\\
点 (拓扑斯的) & IV 6.0\\
已表示预层 & I 1.3.3\\
可表预层 & I 1.4.1\\
预层 (集合) & I 1.2\\
预拓扑 & II 1.3\\
pro-伴随 & I 8.11.5\\
pro-对象 & I 8, I 8.10\\
pro-可表 (函子) & I 8.12\\
pro-可表 (函子) & I 8.10.10\\
投影算子 & I 10.6\\
延空 & III 5.3\\
可作拉回 (箭头、态射) & I 10.7\\
商、严格有效商、普遍商 & I 10.8\\[.7em]
加细 & II 1.1\\
等价关系 & I 10.9\\
局部等价关系 & IV 9.8.3\\
有效等价关系、普遍有效等价关系 & I 10.10\\
限制 (Weil 的) & IV 5.2\\
限制 (层的) & IV 5.4, III 5.3\\[.7em]
分离 (预层) & II 2.1\\
景 & II 1.1.5\\
$\underline U$-景 & II 3.0.2\\
局部赋环景 & IV 13.9\\
清醒 (拓扑空间) & IV 4.2.1\\
不交和 & II 4.5\\
% SOURCE physical page537, printed523, visually checked against page-537.png.
子对象、严格子对象 & I 10.11\\
子拓扑斯 & IV 9.1.1\\
互补子拓扑斯 & IV 9.1.13\\
有补子拓扑斯 & IV 9.1.13\\
闭子拓扑斯 & IV 9.3.5\\
局部闭子拓扑斯 & IV 9.4.9\\
开子拓扑斯 & IV 9.2.3\\
特殊化 (点的) & IV 4.2.2\\
对下降稳定 & IV 8.2\\
严格 ind-可表 & I 8.12.1\\
支集 (群的、群截面的) & IV 9.3.5\\
$r$-结构 & IV 9.8.3\\
支集、余支集 & IV 8.5\\[.7em]
拓扑 & II 1.1\\
$\underline U$-拓扑 & II 3.0.2\\
典范拓扑 & II 2.5\\
离散拓扑 & II 1.1.4\\
粗拓扑与混沌拓扑 & II 1.1.4\\
诱导拓扑 & III 3.1\\
拓扑斯 & IV 1.1\\
$\underline U$-拓扑斯 & IV 1.1\\
赋环拓扑斯 & IV 11.1.1\\
分类拓扑斯 & IV 2.3, 2.4, 2.5\\
等价拓扑斯 & IV 3.4\\
有限拓扑斯 & IV 9.1.12\\
始拓扑斯、终拓扑斯 & IV 2.2\\
刚性拓扑斯 & IV 9.8.2 b)\\[.7em]
宇宙 & I 0\\
宇宙 & I 11.1\\
宇宙 (公理) & I 11.4\\[.7em]
邻域 & IV 6.8\\
\end{longtable}
\endgroup

\chapter*{符号索引}
\addcontentsline{toc}{chapter}{符号索引}
\begingroup
\small
\setlength{\LTleft}{0pt}
\setlength{\LTright}{0pt}
\begin{longtable}{@{}p{.68\linewidth}p{.27\linewidth}@{}}
% SOURCE physical page538, printed524, visually checked against page-538.png.
$\underline a:\widehat C\longrightarrow\widetilde C$ & II 3.4\\
$\underline a F$ & II 3.5\\
$(B,A,f)$ ($f:A\longrightarrow B$ 为函子) & IV 9.5.1\\
$B_G$ & IV 2.3, 4.5\\
$\widetilde C$ & II 2.1, 6.3.3\\
$C,C',D,\ldots$ 范畴 & \\
$\widehat C,\widehat C_{\underline U}$ & I 1.2\\
& II 6.3.3\\
& IV 2.6, 4.6\\
$\widetilde C_A,\widetilde C_{\mathrm{ab}}$ & II 6.3.3\\
$C/F$ & I 3.4.0\\
$\operatorname{Cov}(X)$ & II 1.3\\
$\widetilde C_{\underline U}$ & II 3.6\\
$A^E,E_A,E_{\mathrm{ab}}$ ($E$ 为拓扑斯) & IV 11.1.2\\
$E_{/X}$ & IV 5.1\\
$\underline{\operatorname{Fib}}(E)$ & IV 6.2\\
$\operatorname{Filt}(E)$ & I 9.12\\
$\Gamma(E,X),\Gamma(X)$ & IV 4.3.6\\
$\operatorname{Hom}_A(M,N),\operatorname{Hom}(M,N)$ & IV 12.1\\
$\operatorname{Hom}_E(X,Y),\operatorname{Hom}(X,Y)$ & IV 10.2\\
$\underline{\operatorname{Homtopan}}(E,E'),\underline{\operatorname{Homtoplocan}}(E,E')$ & IV 13.9 c)\\
$\underline{\operatorname{Homtop}}(E,E')$ & IV 3.2, 3.3\\
$h(X),h_{\underline U}(X)$ & I 1.3\\
$i:\widetilde C\hookrightarrow\widehat C$ & \\
$\operatorname{Ind}(C),\operatorname{Ind}_{\underline U}(C),\operatorname{Ind}_{\underline V}(C,\underline U)$ & I 8.2.4.4\\
$\xrightarrow{\sim}$ 同构 & \\
$j^{\mathrm{ab}}_{X!},j^{\mathrm{ens}}_{X!}$ & IV 11.3.2\\
$J(X)$ & II 1.1\\
$j_X:E/X\longrightarrow E$ & IV 5.2\\
$j_{X!},j_X^*,j_{X*}$ & III 11.2.2\\
& IV 5.2.2, 11.2.2, 11.3.1\\
% SOURCE physical page539, printed525, visually checked against page-539.png.
(Karfiness) & IV 9.1.12\\
$L$ (函子) & II 3.1\\
$\operatorname{Lib}$ (妇女解放), 另见 & I 5.8.2, II 6.5\\
$\displaystyle\varprojlim_I$ & I 2.1\\
$\displaystyle\varinjlim_I$ (归纳极限) & \\
$\operatorname{loc}(f)$ & IV 5.5.2\\
$\operatorname{loc}(R),\operatorname{glob}(r)$ & IV 9.8.3 b)\\
$M\otimes_A N$ & IV 12.7\\
$\hookrightarrow$ 单态射或函子 & \\
$\underline{\operatorname{Morsite}}(C',C)$ & IV 4.9.3\\
$\operatorname{Ouv}(X)$ & IV 2.1\\
$\varnothing_{\mathrm{top}}$ & IV 4.4\\
$\Phi_v,\Psi_v,\Xi_v,\Lambda_v$ & IV 10.5\\
$\underline{\operatorname{Points}}(E)$ & IV 4.6.3, 6.1\\
$\underline{\operatorname{Pt}}(E)$ & IV 6.1\\
$\underline{\operatorname{Ptgeom}}(E)$ & IV 13.9 f)\\
$\operatorname{Pro}(C),\operatorname{Pro}_{\underline U}(C),\operatorname{Pro}_{\underline V}(C,\underline U)$ & I 9.10.2\\
$\operatorname{Quot}(U),Q_X$ & IV 9.8.3 a)\\
$\operatorname{Res}(h)$ & IV 11.1.4\\
$(\operatorname{Topfin})$ & IV 9.1.12\\
$\operatorname{Top}(X/r)$ & IV 9.8.3 h)\\
$\operatorname{Top}(f)$ & IV 4.9.1, 4.1.1\\
$\operatorname{Top}(X)$ & IV 4.1, 2.1, 4.2\\
$\operatorname{TOP}(f)$ & IV 4.1.3\\
$\operatorname{TOP}(X)$ & IV 2.5, 4.1\\
$u^*$ & I 5.0\\
$u^*,u_*,u_!$ & I 5.1, 5.9\\
& IV 3\\
$u^{-1},u^*,u_*$ & IV 13.2.3\\
$u_s,u^*$ & III 1.1.1\\
$\underline U$-$\operatorname{Ens}$ & I 0\\
$\underline U,\underline V$ 宇宙 & I 0, 11\\
$(\underline V,\underline U\text{-}\operatorname{Top})$ & IV 3.3\\
$\underline{\operatorname{Vois}}(p)$ & IV 6.8\\
$X_{\complement U}$ & IV 9.3.2\\
$(X'/U)$ & IV 5.2\\
\end{longtable}
\endgroup
